% !TeX root = ../main.tex

\section{Conclusões}
\label{sec:conclusoes}

O estudo modelou o cronograma de projeto de um empreendimento de grande
porte como um RCPSP e mediu quanto o prazo pode cair apenas com a escolha
dos dias de início, mantendo as mesmas atividades, durações, vínculos e
equipes. Respeitando as datas em que as disciplinas foram liberadas no
cronograma praticado, o prazo ótimo é de 285 dias, contra 310 do
praticado: uma redução de 25 dias (8,1\%). Se todas as disciplinas
pudessem começar desde o início, o ótimo seria de 187 dias, 123 a menos.
A diferença entre os dois valores é a espera das disciplinas pela base de
arquitetura, que, na realidade, conduz o processo de projeto.

Nos dois casos, o prazo é definido pela sequência de emissões e análises,
e não pela capacidade das equipes: uma frente a mais em todas as equipes
não muda o prazo. A alavanca é o tempo de emissão do projetista nas
atividades da cadeia crítica, que na versão realista é a do anteprojeto de
arquitetura e caixilhos. Emitir 10\% mais rápido reduz 15 dias, enquanto
reduzir à metade o tempo de análise do cliente reduz 2. Uma redução fora
da cadeia crítica não altera o prazo. Daí a principal implicação prática:
padronização e automação dão resultado quando são dirigidas às atividades
que a análise de folgas aponta, e não quando são aplicadas de forma geral.

Do ponto de vista do método, a formulação indexada no tempo forneceu um
limite inferior justo já na relaxação linear, mas o solver não encontrou
cronogramas bons para 243 atividades em 36 minutos. A combinação de
caminho crítico e esquema serial resolveu a instância completa em menos
de um segundo, com prova de otimalidade pela igualdade entre o limite
inferior e o prazo encontrado. Na instância pequena, em que a capacidade
limita o prazo, foi o solver que provou o ótimo. Os dois métodos se
completam.

As principais limitações são a ausência, na base, do bloco do estudo
preliminar, substituído por datas de liberação, as capacidades estimadas
pelo pico do praticado e o tratamento das durações e do número de revisões
como dados. Como trabalhos futuros, propõe-se incluir o estudo preliminar
no modelo, para medir o efeito de liberar as disciplinas por pacotes,
tratar o número de revisões como decisão, e usar o cronograma heurístico como solução inicial do solver em
instâncias em que a capacidade esteja ativa.
